%========================== APÊNDICES ==========================
% Correção de sumário: os apêndices entram como SUBSEÇÕES não numeradas e a entrada
% no sumário é adicionada como TEXTO ("APÊNDICE X — ..."), não como número. Assim o
% rótulo longo não estoura a caixa de número do \numberline (causa da sobreposição).
\newpage
\phantomsection
\addcontentsline{toc}{section}{APÊNDICES}
\section*{APÊNDICES}

\newcounter{apcount}
\newcommand{\apendice}[1]{%
  \par\bigskip
  \phantomsection
  \refstepcounter{apcount}%
  \subsection*{APÊNDICE \Alph{apcount}\ --- #1}%
  \addcontentsline{toc}{subsection}{APÊNDICE \Alph{apcount}\ --- #1}%
}

%------------------------------------------------------------------
\apendice{Roteiro de tarefas do teste de usabilidade (SUS)}\label{ap:sus}

Cada participante executa, sem treinamento prévio, as cinco tarefas representativas
abaixo; ao final, responde ao questionário \textit{System Usability Scale} (SUS),
de dez itens em escala Likert de 1 (discordo totalmente) a 5 (concordo totalmente).

\begin{enumerate}
  \item Comissionar (adicionar) um dispositivo à casa a partir da descoberta na rede.
  \item Ligar e desligar um aparelho por \textbf{comando de voz} em português.
  \item Criar uma \textbf{automação por horário} (rotina) para um dispositivo.
  \item Consultar o \textbf{painel de energia} e identificar o consumo de um aparelho.
  \item Recuperar-se de uma \textbf{falha simulada} (dispositivo \textit{offline}).
\end{enumerate}

Os dez itens do SUS e as tarefas acima estão implementados no instrumento de coleta
descrito no Apêndice~\ref{ap:survey}. A pontuação e a discussão constam da
Seção~\ref{sec:res-sus} (\textit{a coletar}).

%------------------------------------------------------------------
\apendice{Corpus de comandos de voz}\label{ap:corpus}

Protocolo do corpus controlado para o cálculo da taxa de erro de palavra (WER) e da
matriz de confusão de intenções: no mínimo \textbf{50 enunciados} distintos, gravados
por \textbf{pelo menos três falantes} (variando sotaque e idade), cada um com
transcrição de referência (\textit{ground truth}) e intenção-verdade rotulada
manualmente. As categorias de intenção são: \emph{ligar}, \emph{desligar},
\emph{alternar}, \emph{ajustar brilho} e \emph{ajustar cor}. Lista completa e áudios
\textit{a coletar} (ver Seção~\ref{sec:res-voz}).

%------------------------------------------------------------------
\apendice{Instrumento de pesquisa de usuários}\label{ap:survey}

Para validar as premissas do trabalho com usuários reais de diferentes idades e
regiões (Recife/PE, São Paulo, Rio de Janeiro e participantes na Itália/Sicília), foi
elaborado um questionário aplicado por formulário eletrônico (Google Forms), anônimo e
em conformidade com a LGPD. A coleta caracteriza uma avaliação \textbf{direcional}
(amostra de conveniência, $n = 8$--$15$), não uma inferência estatística. As respostas
estavam \textbf{em coleta} no fechamento desta versão; a consolidação e os gráficos
constam da Seção~\ref{sec:res-survey}.

\vspace{0.4em}
\noindent\textbf{Bloco 1 --- Perfil:} faixa etária (18--24 / 25--34 / 35--44 / 45--59 /
60+); cidade ou região; escolaridade; familiaridade com tecnologia (1--5); faixa de
renda familiar (opcional).

\noindent\textbf{Bloco 2 --- Situação atual:} possui algum dispositivo de casa
inteligente? (sim/não); quais; quem os configurou; frequência de uso.

\noindent\textbf{Bloco 3 --- Intenção e barreiras:} automatizaria a casa se fosse simples
e barato? (1--5); \emph{maiores dificuldades percebidas} (preço do \textit{hardware};
complexidade de instalar/parear; depender de nuvem; privacidade; falta de conhecimento
técnico; \textit{app} do fabricante confuso; medo de indisponibilizar a casa); quanto
pagaria por uma solução completa; preferência por operação \textit{offline}/local;
preocupação com privacidade (1--5).

\noindent\textbf{Bloco 4 --- Avaliação do DOMUS (para quem viu a demonstração):}
os dez itens do SUS (Apêndice~\ref{ap:sus}); ``conseguiria controlar por voz sem o
\textit{app} do fabricante?''; comentário aberto.

\vspace{0.4em}
\noindent Formulário público (instrumento aplicado):
\newline\url{https://docs.google.com/forms/d/e/1FAIpQLSd3DUqWC2W06tpbYuAnqMaZheARLSYvXE7C5gBiCwN0HuOzag/viewform}

%------------------------------------------------------------------
\apendice{Lista de materiais e custos (BOM)}\label{ap:bom}

Planilha datada, com preços de varejo referenciados (nota fiscal ou \textit{link}),
em dois cenários: \emph{incremental} (computador doméstico reaproveitado como
\textit{hub}, custo marginal zero) e \emph{completo} (mini-PC/Raspberry~Pi dedicado).
Valores \textit{a coletar} (ver Seção~\ref{sec:res-custo}, Tabela~\ref{tab:bom}).

%------------------------------------------------------------------
\apendice{Repositório e artefatos}\label{ap:repo}

Código-fonte, histórico e instruções de reprodução:
\newline\url{https://github.com/julianosfreitas/DOMUS_HOUSE} (\textit{branch} \texttt{main}).
O \textit{commit} de referência da defesa e o roteiro de execução (\emph{back end}
NestJS + PWA Next.js + PostgreSQL, com modo \textsc{Mock} que dispensa \textit{hardware})
constam do \texttt{README.md} do repositório.
